\section{Mathematical Tools and Notation}

This appendix collects all notation and standard mathematical results used in the main text
and in Appendix A, together with the connections between ``advanced theory'' and the subject
matter. All quoted results are attributed; referenced (non-original) content does not
participate in our grading scheme.

\subsection{Notation Table}

\begin{table}[htbp]
\centering
\caption{Principal notation\label{tab:notation}}
\begin{tabular}{ll}
\toprule
Notation & Meaning \\
\midrule
$\Sigma$ & finite vocabulary (token vocabulary), $\abs{\Sigma} < \infty$ \\
$\Sigma^{*}$ & set of finite strings; $\Sigma^{\le C}$ the set of strings of length at most $C$ \\
$D$ & static pretraining corpus ($\Pi_1$) \\
$A$ & training algorithm, $\theta = A(D)$ \\
$\theta \in \Theta \subseteq \R^{P}$ & model parameters ($P$ the number of parameters) \\
$p_\theta(y_{1:T}\mid c)$ & autoregressive conditional probability (Eq.~\eqref{eq:autoregressive}) \\
$c, c_{\le t}$ & context; context at time $t$ \\
$C$ & context window length ($\Pi_3$) \\
$\Phi_\theta: \Sigma^{*} \to \R^{d}$ & forward function (Definition \ref{def:forward}) \\
$H_c = \Phi_\theta(c)$ & hidden state (residual-stream activation) \\
$R$ & self-report (Definition \ref{def:report}) \\
$O^{*}: \Theta \times \Sigma^{*} \to \Sigma^{*}$ & ground-truth introspection mapping (Definition \ref{def:oracle}) \\
$\ell(\cdot,\cdot)$ & report distortion measure (Definition \ref{def:metrics}) \\
$\mathcal{R}_{\mathrm{int}}$ & introspective risk (Eq.~\eqref{eq:int-risk}) \\
$I(X;Y)$, $I(X;Y\mid Z)$ & (conditional) mutual information \\
$H(X)$, $H(X\mid Y)$ & (conditional) entropy \\
$\TV{}{P,Q}$ & total variation distance \\
$\doop(\cdot)$ & Pearl do-operator\cite{pearl-2009} \\
$X \indep Y \mid Z$ & $X, Y$ conditionally independent given $Z$ \\
$\mathrm{A}_1$--$\mathrm{A}_5$ & the five introspective criteria (Definition \ref{def:criterion}) \\
$\Pi_1$--$\Pi_5$ & the paradigm postulates (Assumption \ref{asm:paradigm}) \\
$\mathrm{CL}(\mathcal{I})$ & closed-loop effect measure (Definition \ref{def:cl}) \\
\bottomrule
\end{tabular}
\end{table}

\subsection{Information-Theoretic Tools}

\begin{lemma}[Maximum likelihood approximates the data distribution]\label{lem:mle}
Let $\theta^{*} \in \arg\min_\theta \E_{x \sim D}[-\log p_\theta(x)]$. If the model family
$\{p_\theta\}$ is sufficiently expressive and the data are sufficient, then $p_{\theta^{*}}$
minimizes $\KL{P_D}{p_\theta}$ ($P_D$ the empirical distribution of the corpus). In general,
for any query context $c$,
\begin{equation}
p_{\theta^{*}}(\cdot \mid c) \;\approx\; P_D(\cdot \mid c),
\label{eq:mle}
\end{equation}
i.e., the model's conditional output distribution approximates the \textbf{corpus-text}
conditional distribution.
\end{lemma}

\begin{proof}[Note]
This is the standard property of maximum likelihood estimation (KL-divergence minimization);
see \cite{cover-thomas-2006} (Chapters 2 and 11) and \cite{shannon-1948}. Eq.~\eqref{eq:mle}
is the formal basis of the statement in Theorem 4a that ``meta-statements are functions of
corpus statistics''. $\square$
\end{proof}

\begin{lemma}[Data Processing Inequality (DPI)]\label{lem:dpi}
For a Markov chain $X \to Y \to Z$, $I(X;Z) \le I(X;Y)$.
For any function $f$, $I(X; f(Y)) \le I(X;Y)$. \cite{cover-thomas-2006}
\end{lemma}

\begin{remark}[Role of the DPI]
Theorem 1a is essentially a ``zero-channel'' statement: $R$ and $H$ are conditionally
independent given $(\theta,c)$, so any information channel from $H$ to an observer through
$R$ has zero capacity, a degenerate case of the DPI. The intuition of Theorem 5c that
``probe readout outperforms text report'' can also be viewed as a DPI variant: the report $R$
is a function of $H$ through $p_\theta$ (the text bottleneck), and its information content is
limited by the bottleneck.
\end{remark}

\subsection{Computational Learning Theory Tools}

\begin{restate}[No-Free-Lunch (NFL)]
Let $\mathcal{X}$ be a finite domain ($\abs{\mathcal{X}} = N$), labels
$\mathcal{Y} = \{0,1\}$, and $D_f$ the uniform distribution on $\mathcal{X}$ (labels given
by the function $f: \mathcal{X} \to \mathcal{Y}$). For any learning algorithm $A$ and any
number of training samples $m \le N/2$,
\begin{equation}
\E_{S \sim (D_f)^m,\, f \sim \mathrm{Unif}(\mathcal{Y}^{\mathcal{X}})}\!\!\left[\, L_{D_f}(A(S)) \,\right] \;\ge\; \frac12 - \frac{m}{2N},
\label{eq:nfl}
\end{equation}
where $L_{D_f}$ is the 0-1 generalization loss. In particular, for any algorithm there exists
a label function $f$ with expected error $\ge \tfrac12 - \tfrac{m}{2N}$; when $m = 0$ (no
samples), this lower bound is $\tfrac12$.
\end{restate}

\begin{proof}[Proof sketch]
For a fixed training set $S$ ($m$ points), the labels of the $N - m$ unlabeled points admit
$2^{N-m}$ equally likely assignments; the algorithm's predictions on the unlabeled points
agree in expectation with random guessing. Averaging over $S$ and the concept uniformly and
using symmetry yields Eq.~\eqref{eq:nfl}. For the complete proof see \cite{shalev-shwartz-2014}
(Theorem 5.1) and \cite{wolpert-macready-1997}. $\square$
\end{proof}

\begin{remark}[The sample requirement of PAC learning]\label{thm:pac}
(Valiant\cite{valiant-1984}; form given in \cite{shalev-shwartz-2014}) PAC learning
presupposes \textbf{labeled samples}: if the training set is empty, any generalization
guarantee of a learning algorithm over target functions degenerates to the NFL lower bound
(Eq.~\eqref{eq:nfl}), one of the direct bases of Theorem 1b.
\end{remark}

\begin{remark}[The PAC-Bayes perspective]
The PAC-Bayes theorem\cite{mcallester-1999} bounds the generalization error by the empirical
error plus a term of the order of the square root of the ``prior $\pi$ to posterior $A(S)$ KL
divergence''
($\hat{L}_S(A) + O\!\left(\sqrt{(\KL{\pi}{A(S)} + \ln \tfrac{\sqrt{m}}{\delta})/m}\right)$,
one of McAllester's forms). In the present context, the ``samples'' of the
introspective task satisfy $S = \emptyset$ ($\Pi_5$), the empirical error $\hat{L} = 0$, and
the prior--posterior KL term carries no information, and PAC-Bayes likewise degenerates to no
guarantee. Theorem 1b is the constructive formulation of this observation.
\end{remark}

\subsection{Goodhart's Law and Objective Mismatch}

\begin{remark}[Taxonomy of Goodhart effects]\label{thm:goodhart-tax}
Manheim--Garrabrant\cite{manheim-garrabrant-2018-goodhart} classify Goodhart effects into
regressional, causal, adversarial, and extremal types: the \textbf{regressional} type refers
to the situation where ``when the proxy metric $M$ and the true objective $T$ are only
statistically correlated, extreme optimization of $M$ drives the expectation of $T$ toward
its prior mean rather than its optimum (and actively deteriorates $T$ when $M$ and $T$ are
negatively correlated)''. Theorem 2 is a strict instantiation of the
regressional type.
\end{remark}

\begin{restate}[Characterization of reward hacking (Skalse et al.)]
(\cite{skalse-2022-reward-hacking}) For a proxy reward $\widetilde{R}$ and a true reward
$R$: if $\widetilde{R}$ does not imply $R$ (there exist state pairs whose order relations
differ), then, over sufficiently expressive policy classes, policies maximizing
$\widetilde{R}$ can arbitrarily harm the expected value of $R$. A formal version requires
expressibility and state-coverage conditions.
\end{restate}

\begin{remark}[Relation to Theorem 2]
Theorem 2 is an \textbf{instantiation} of the two results above: the proxy metric
$M = -\ell(\theta)$ (text likelihood), the true objective $T = -\mathcal{R}_{\mathrm{int}}(\theta)$
(introspective fidelity). The key difference: the present work \textbf{proves} their statistical
relation (Eq.~\eqref{eq:oracle-indep}, $O^{*} \indep D \mid \theta$), making the ``mismatch''
a structural fact of the paradigm rather than an empirical observation.
\end{remark}

\subsection{Causal Inference Tools}

\begin{definition}[Do-operator]\label{def:do}
(Pearl\cite{pearl-2009}) $P(Y \mid \doop(X = x))$ denotes the interventional distribution of
``forcibly setting $X$ to $x$'': delete all edges pointing into $X$ in the causal graph, set
$X = x$, and compute the distribution of $Y$. The interventional distribution generally
differs from the observational distribution $P(Y \mid X = x)$.
\end{definition}

\begin{lemma}[Interventions on deterministic functions]
If $Y = f(X)$ is a deterministic function, then for any $x \ne x'$,
$P(Y \mid \doop(X = x)) = \delta_{f(x)}$, and
$\TV{}{\delta_{f(x)}, \delta_{f(x')}} = \mathbf{1}\{f(x) \ne f(x')\}$.
\end{lemma}

\begin{remark}[Use of interventions]
All arguments of Theorem 3 rest on the characterization of the ``legal intervention set'':
the only quantity intervenable within paradigm $\Pi$ is the context text ($\Pi_2, \Pi_3$).
Intervention on the hidden state $H$ (the paper's concept injection) is an
\textbf{out-of-paradigm} intervention, executed by an external experimenter, whose very
existence corroborates the ``no first-person channel'' conclusion of Theorem 5 (externally
intervenable = third-party accessible).
\end{remark}

\subsection{Algorithmic Information Theory and Recursion Theory Background}

The following classical results are \textbf{not directly invoked} in the proofs of this
report, but constitute the ``advanced-theory'' background and are relevant to the subject:

\begin{enumerate}
\item \textbf{Kolmogorov complexity}\cite{kolmogorov-1965}: the complexity of a string $x$ is
$K(x) = \min\{|p| : U(p) = x\}$. Background connection: the complexity of the model's
self-model $g_\theta$ does not exceed $K(\theta) + K(\text{query})$; at inference,
$K(g_{\theta(t)})$ is constant (the algorithmic-information-theoretic counterpart of Theorem
6a). If self-knowledge required an ``incompressible description of oneself'', its complexity
lower bound would grow linearly with self-history, so a static system cannot satisfy this
(Theorem 6b).
\item \textbf{Rice's theorem}\cite{rice-1953}: nontrivial semantic properties of programs are
undecidable. Background connection: if ``whether a model introspects'' is interpreted as a
semantic property of the function the model computes, then no universal adjudication
procedure exists; this supports the cautious stance of the paper and of the philosophical
literature that ``introspective capacity cannot be directly adjudicated from behavior'', but
our theorems take the stronger route (structural impossibility rather than
undecidability).
\item \textbf{Identification in the limit (Gold)}\cite{gold-1967}: language identification in
the limit requires an \textbf{infinite evidence stream}. Background connection: the runtime
evidence stream of paradigm $\Pi$ is the context ($\le C$ tokens), finite and bounded, so
under the Gold framework one cannot identify ``one's own language'' (self-model), complementing
Theorem 6b.
\item \textbf{The free-energy principle (Friston)}\cite{friston-2010}: biological agents
continuously minimize variational free energy about their own states. Background connection:
human introspection can be formalized as ``online variational inference about one's own
state'', whose update law depends on \textbf{continuous sensory and proprioceptive streams};
paradigm $\Pi$ has no such streams (Theorem 1a) and no runtime free-energy-minimization
mechanism (Theorem 3); the free-energy principle provides a dynamical-systems basis for
``human introspection requires online self-inference'' (grade-D background).
\end{enumerate}

\clearpage
